
%\section{Pipeline Dynamics}
%\label{app:pipeline-dynamics}

\subsubsection{Purpose and Scope}
The discussion in the Appendix till now formalizes
\emph{static} correctness of individual FRFP pipelines, Epistemic Dynamics
analyzes what happens when such pipelines are iterated by the same
human--AI system over many turns. The focus is on:
\begin{enumerate}
  \item persistence and degradation of tacit context;
  \item explicit-only hallucinations and their containment;
  \item confidence inflation and requirement escalation;
  \item feedback-coupled degradation and accelerated hallucination;
  \item stochastic and finite-horizon pipelines.
\end{enumerate}
No collective or institutional mechanisms are assumed in this section;
these appear in latter.

Throughout, we fix the Epistemic structure and work
with the underlying tacit context space \(T\) and explicit artifact
space \(E\).

%======================================================================
\subsubsection{Tacit Context Space}
%======================================================================

Each human agent possesses a tacit context space \(T\), equipped with a
preorder \(\preceq\).

\begin{definition}[Tacit context preorder]
The binary relation \(\preceq\) on \(T\) is read as:
\[
t_1 \preceq t_2
\quad\text{``\(t_1\) is epistemically no stronger than \(t_2\)''}.
\]
We write \(t_1 \prec t_2\) if \(t_1 \preceq t_2\) and \(t_2 \not\preceq t_1\).
\end{definition}

No algebraic structure on \(T\) is assumed beyond this preorder. The
categorical tacit structure  (one-object category,
primitives \(\TE,\HFD\), correctness quotient \(\gamma\)) is still in
force, but for dynamic analysis we work primarily with the underlying
preordered set \((T,\preceq)\).

%======================================================================
\subsubsection{Context Degradation}
%======================================================================

Tacit context degrades over time due to forgetting, misplaced trust,
habituation, and other long-run effects.

\begin{definition}[Context degradation]
A \emph{context degradation} operator is a function
\[
\delta : T \to T
\]
modeling loss of grounding across turns. We assume:
\begin{enumerate}
  \item \textbf{Monotonicity}:
        \(\delta(t) \preceq t\) for all \(t \in T\).
  \item \textbf{Strict decay somewhere}:
        there exists \(t \in T\) with \(\delta(t) \prec t\).
  \item \textbf{Non-invertibility}:
        \(\delta\) has no right inverse; i.e., there is no function
        \(f:T\to T\) such that \(\delta \circ f = \mathrm{id}_T\).
\end{enumerate}
We write \(\delta^n\) for the \(n\)-fold iterate of \(\delta\).
\end{definition}

Intuitively, \(\delta\) expresses the part of tacit dynamics that is
\emph{purely degrading}. In this appendix we deliberately ignore
repair/refresh operations (e.g., reading textbooks, taking breaks) to
obtain worst-case inevitability results. Repair can be modeled by
separate operators \(\rho:T\to T\), but these are left for future work.

%======================================================================
\subsubsection{Dynamic Interpretation of Pipelines}
%======================================================================

We now interpret complete FRFP pipelines as transformations of tacit
context.

Let \(E_{\mathrm{sem}}\) be the explicit semantic category and \(\RB^\sharp : E_{\mathrm{sem}}\to T\) the semantic
backflow (Section~A.0.5). For each explicit pipeline \(e\in
E_{\mathrm{sem}}\), the Epistemic Space suffix \(\TE,\HFD\) yields a tacit
transformer
\[
t \longmapsto \HFD \circ \TE \circ \RB^\sharp(e,t),
\]
where we allow the tacit state \(t\) to modulate interpretation
(e.g., via background knowledge).

\begin{definition}[Pipeline transformer]
A \emph{complete FRFP pipeline} \(P\) at a given turn consists of:
\begin{itemize}
  \item an explicit artifact \(e \in E_{\mathrm{sem}}\) produced by the
        AI system via explicit computation and navigation;
  \item the canonical tacit suffix \(\RB \to \TE \to \HFD\);
  \item subsequent degradation \(\delta\).
\end{itemize}
It induces a tacit transformer
\[
\llbracket P \rrbracket_\delta : T \to T,\qquad
\llbracket P \rrbracket_\delta(t)
 :=
 \delta\bigl( \HFD\circ\TE\circ\RB^\sharp(e,t)\bigr).
\]
\end{definition}

\begin{definition}[Run]
A \emph{run} (or \emph{pipeline execution}) is a sequence
\[
t_0, t_1, t_2,\dots
\]
of tacit states such that there exist pipelines
\(P_0, P_1, P_2,\dots\) with
\[
t_{n+1}
  = \llbracket P_n \rrbracket_\delta(t_n)
\quad\text{for all } n \geq 0.
\]
\end{definition}

We will often also track the explicit artifacts \(e_n\) and credences
\(c_n\) associated with each \(P_n\); this is made precise in
Sections~\ref{sec:B7-confidence} and~\ref{sec:B15-stochastic}.

%======================================================================
\subsubsection{Requirements and Hallucination}
%======================================================================

We formalize hallucination in terms of a requirement map on explicit
artifacts.

\begin{definition}[Requirement map]
A \emph{requirement map} is a function
\[
\Req : E \times \mathbb{R}_{\ge 0} \to T
\]
assigning to each explicit artifact \(e \in E\) and tacit credence
\(c\ge 0\) a required tacit grounding \(\Req(e,c)\). We require:
\begin{enumerate}
  \item \textbf{Monotonicity in credence}:
        if \(c_1 \le c_2\) then
        \(\Req(e,c_1) \preceq \Req(e,c_2)\);
  \item \textbf{Nontriviality}:
        there exist artifacts whose requirements strictly exceed others
        in the preorder.
\end{enumerate}
When credence is suppressed, we write \(\Req(e)\) for a baseline
requirement (e.g., at default credence).
\end{definition}

\begin{definition}[Hallucination]
Let \(t\in T\) be the tacit state at evaluation time, \(e\in E\) an
explicit artifact, and \(c\ge 0\) the tacit credence assigned to
\(e\). A \emph{hallucination} occurs at this step if
\[
t \prec \Req(e,c).
\]
We say that the artifact \(e\) is \emph{hallucinated} (relative to
\((t,c)\)) in this case.
\end{definition}

\begin{remark}[Structural nature of hallucination]
In this formalism, hallucination is an \emph{implementation-agnostic}
relation between tacit state \(t\), requirement \(\Req(e,c)\), and
explicit artifact \(e\). It does not depend on the model's internal
computation or intent. Concrete architectures (e.g., LLMs) shape
\emph{how often} trajectories enter states with \(t \prec
\Req(e,c)\), but do not appear in the definition itself.
\end{remark}

%======================================================================
\subsubsection{Hallucination Inevitability}
%======================================================================

We now state a baseline hallucination inevitability theorem for
deterministic degradation dynamics.

To capture recurring nontrivial demands, we restrict to runs that keep
encountering artifacts whose requirements exceed a fixed lower bound.

\begin{definition}[Nontrivial requirement floor]
A requirement floor is a tacit state \(r\in T\). A run
\((t_n,e_n,c_n)_{n\ge 0}\) has \emph{recurring nontrivial demands}
relative to \(r\) if there exists an infinite subsequence
\((n_k)_{k\ge 0}\) such that
\[
\Req(e_{n_k},c_{n_k}) \succeq r
\quad\text{for all }k.
\]
\end{definition}

We also assume that degradation eventually pushes tacit states below
this floor.

\begin{assumption}[Baseline degradation below the requirement floor]
\label{assump:baseline-degradation}
There exists \(r\in T\) such that for every initial tacit state
\(t_0\in T\), there exists \(N\in\mathbb{N}\) with
\[
\delta^N(t_0) \prec r.
\]
\end{assumption}

This is a strengthening of ``strict decay somewhere'': it requires that
every trajectory ultimately falls below the floor \(r\).

\begin{theorem}[Hallucination inevitability]
\label{thm:baseline-inevitability}
Assume:
\begin{enumerate}
  \item context degradation \(\delta\) satisfies monotonicity,
        non-invertibility, and
        Assumption~\ref{assump:baseline-degradation};
  \item requirements admit a nontrivial floor \(r\) and are monotone in
        credence;
  \item a run \((t_n,e_n,c_n)_{n\ge 0}\) has recurring nontrivial
        demands relative to \(r\).
\end{enumerate}
Then every infinite run produces at least one hallucination; i.e., there
exists \(n\) with
\[
t_n \prec \Req(e_n,c_n).
\]
\end{theorem}

\begin{proof}[Proof sketch]
By Assumption~\ref{assump:baseline-degradation}, for any initial
\(t_0\) there exists \(N\) with \(\delta^N(t_0) \prec r\). Consider the
given run and the subsequence \((n_k)\) along which
\(\Req(e_{n_k},c_{n_k})\succeq r\). For sufficiently large \(k\),
monotonicity of \(\delta\) and recurrence of nontrivial demands imply
that
\[
t_{n_k} \preceq \delta^{m}(t_0) \prec r \preceq \Req(e_{n_k},c_{n_k})
\]
for some \(m\), so \(t_{n_k} \prec \Req(e_{n_k},c_{n_k})\), i.e., a
hallucination. A detailed proof can be given by induction along the
run, but the key idea is that degradation eventually drives tacit
grounding below a floor that requirements exceed infinitely often.
\end{proof}

\begin{remark}
The theorem applies to \emph{nontrivial, open-ended} task streams with
recurrent demands above a fixed floor. Trivial avoidance regimes (e.g.,
restricting to a tiny class of always-easy tasks or always answering
``I don't know'') fall outside the intended scope.
\end{remark}

%======================================================================
\subsubsection{Explicit-Only Hallucinations}
%======================================================================
We now isolate explicit artifacts that are produced entirely within explicit reasoning, but which can be
hallucinated relative to some tacit state encountered along a run.

\begin{definition}[Explicit-only hallucination]
An \emph{explicit-only hallucination} is an explicit artifact
\(e \in E\) such that:
\begin{enumerate}
  \item \(e\) is well-typed and syntactically valid as an FRFP explicit
        object;
  \item \(e\) is surface-plausible under an explicit plausibility
        functional (Section~\ref{sec:B7-confidence});
  \item there exists a tacit state \(t\) and credence \(c\) encountered
        in some run such that
        \(t \prec \Req(e,c)\).
\end{enumerate}
\end{definition}

\begin{proposition}[Containment but not elimination]
Nothing in Epistemic Space or Epistemic Algebra forbids the production of explicit-only
hallucinations. Eliminating them entirely would require either:
\begin{itemize}
  \item an explicit correctness oracle, or
  \item a simulation of tacit evaluation inside explicit space,
\end{itemize}
both of which would violate Explicit--Tacit Separation (ETS) and
Human-Exclusive Grounding (HEG).

The explicit--tacit separation guarantees \emph{containment}, not
\emph{elimination}: explicit-only hallucinations cannot become
correctness judgments without passing through \(\RB\), but they may
still be produced arbitrarily often within explicit space.
\end{proposition}

%======================================================================
\subsubsection{Confidence Inflation}
\label{sec:B7-confidence}
%======================================================================

We model how plausible explicit artifacts raise tacit credence even
when they do not increase grounding.

\begin{definition}[Explicit plausibility]
An \emph{explicit plausibility functional} is a map
\[
\pi : E \to [0,1],
\]
measuring surface coherence, fluency, or apparent authority of explicit
artifacts.
\end{definition}

We extend tacit state with a credence component.

\begin{definition}[Extended tacit state]
Let
\[
T^e := T \times \mathbb{R}_{\ge 0}
\]
be the space of extended tacit states \((t,c)\), where \(t\in T\) is
grounding and \(c\) is tacit credence (how settled or reliable the
agent treats the current artifact).
\end{definition}

\begin{definition}[Confidence inflation]
A \emph{confidence inflation operator} is a map
\[
\Phi : (T^e \times E) \to T^e,\qquad
\Phi\bigl( (t,c), e \bigr)
 :=
 (t,\; c + \varphi(\pi(e))),
\]
where \(\varphi:[0,1]\to \mathbb{R}_{\ge 0}\) is monotone increasing.

We require that \(\Phi\) does not increase tacit grounding:
\[
\Phi\bigl((t,c),e\bigr) = (t',c') \quad\Rightarrow\quad t' = t.
\]
\end{definition}

Thus confidence inflation changes only the agent's interpretive stance
(higher credence), not the underlying grounding \(t\).

\begin{remark}
Confidence inflation captures the structural asymmetry between
``feeling confident'' and ``being grounded''. Systems optimized for
producing well-formed, plausible outputs (e.g., LLMs) inevitably
exhibit this asymmetry: they can increase credence by producing
plausible artifacts without necessarily increasing grounding.
\end{remark}

%======================================================================
\subsubsection{Requirement Escalation Under Plausibility}
%======================================================================

Credence affects what requirements are appropriate: asserting
high-credence claims demands stronger grounding.

\begin{definition}[Requirement escalation]
The requirement map \(\Req : E \times \mathbb{R}_{\ge 0} \to T\) is
\emph{escalatory} if it is monotone in credence, i.e.,
\[
c_1 \le c_2 \quad\Rightarrow\quad
\Req(e,c_1) \preceq \Req(e,c_2),
\]
for all \(e\in E\). Repeated exposure to highly plausible artifacts
(large \(\pi(e)\)) can then raise requirements even if the explicit
content \(e\) is unchanged.
\end{definition}

Combining confidence inflation and requirement escalation, we see that
the same explicit artifact can acquire higher requirements over time as
agents become more confident in it.

%======================================================================
\subsubsection{Feedback-Coupled Degradation}
%======================================================================

We now allow degradation to depend on credence, introducing a feedback
loop: overconfidence reduces skepticism and checking, accelerating loss
of grounding.

\begin{definition}[Feedback-coupled degradation]
A \emph{feedback-coupled degradation} operator is a function
\[
\delta^\star : T^e \to T,\qquad
\delta^\star(t,c) = \delta_c(t),
\]
where \(\delta_c : T \to T\) is a family of degradation operators
indexed by credence \(c\), satisfying:
\begin{enumerate}
  \item \textbf{Credence monotonicity}:
        for fixed \(t\), if \(c_1 \le c_2\) then
        \(\delta^\star(t,c_2) \preceq \delta^\star(t,c_1)\);
        i.e., higher confidence weakens effective grounding.
  \item \textbf{Consistency with baseline degradation}:
        \(\delta^\star(t,0) = \delta(t)\) for all \(t\).
\end{enumerate}
\end{definition}

\begin{definition}[Extended run with feedback]
An \emph{extended run} is a sequence
\[
(t_0,c_0,e_0), (t_1,c_1,e_1), \dots
\]
such that, at each step \(n\),
\begin{align*}
(t_n^+, c_n^+) &= \Phi\bigl((t_n,c_n), e_n\bigr)
      &&\text{(confidence inflation)},\\
t_{n+1} &= \delta^\star(t_n^+, c_n^+)
      &&\text{(feedback-coupled degradation)},\\
c_{n+1} &\text{ chosen by tacit update (e.g., reset or carry-over)}.
\end{align*}
\end{definition}

The exact rule for updating \(c_{n+1}\) can be left abstract; all we
require is that confidence inflation can push credence arbitrarily high
along some runs.

%======================================================================
\subsubsection{Rate of Context Loss}
%======================================================================}

We can analyze the rate at which tacit grounding is lost under
feedback-coupled degradation.

\begin{definition}[Grounding measure and loss rate]
Let \(\mu : T \to \mathbb{R}_{\ge 0}\) be any monotone measure of tacit
grounding (used only for analysis). For an extended run
\((t_n,c_n,e_n)\), define the instantaneous loss rate at turn \(n\) by
\[
\rho_n := \mu(t_n) - \mu\bigl(\delta^\star(t_n,c_n)\bigr).
\]
\end{definition}

\begin{proposition}[Loss rate increases with credence]
Under the credence-monotonicity assumption for \(\delta^\star\) and
monotonicity of \(\mu\), the loss rate \(\rho_n\) is non-decreasing in
\(c_n\). In particular, higher credence does not slow context loss and
may accelerate it.
\end{proposition}

\begin{proof}
If \(c_1 \le c_2\), then
\(\delta^\star(t,c_2) \preceq \delta^\star(t,c_1)\) and
monotonicity of \(\mu\) yields
\(\mu(\delta^\star(t,c_2)) \le \mu(\delta^\star(t,c_1))\). Thus
\[
\mu(t) - \mu(\delta^\star(t,c_2))
\;\ge\;
\mu(t) - \mu(\delta^\star(t,c_1)),
\]
so \(\rho_n\) increases with \(c_n\).
\end{proof}

%======================================================================
\subsubsection{Accelerated Hallucination}
%======================================================================

We now combine confidence inflation, requirement escalation, and
feedback-coupled degradation.

\begin{definition}[Feedback-coupled update]
Along an extended run, the joint update
\[
(t_{n+1}, c_{n+1})
 =
\Bigl(
  \delta^\star\bigl( t_n, c_n + \varphi(\pi(e_n))\bigr),\;
  c_{n+1}
\Bigr)
\]
captures the combined effect of:
\begin{itemize}
  \item confidence inflation
        \((t_n, c_n) \mapsto (t_n, c_n + \varphi(\pi(e_n)))\);
  \item degradation \(\delta^\star\), which is more severe at higher
        credence.
\end{itemize}
Grounding weakens faster while requirements grow stronger (via
\(\Req(e_n,c_n)\)).
\end{definition}

This creates a widening gap between available grounding \(t_n\) and
demanded grounding \(\Req(e_n,c_n)\), leading to accelerated
hallucination.

%======================================================================
\subsubsection{Accelerated Hallucination Inevitability}
%======================================================================

We can now state a strengthened inevitability theorem in the presence
of feedback.

\begin{assumption}[Unbounded credence under plausibility]
There exists \(\bar{\pi} \in (0,1]\) such that for any initial credence
\(c_0\) and any infinite sequence of artifacts with
\(\pi(e_n) \ge \bar{\pi}\) infinitely often, the accumulated credence
\[
c_n^{+} := c_0 + \sum_{k=0}^{n-1} \varphi(\pi(e_k))
\]
tends to \(+\infty\) along a subsequence.
\end{assumption}

\begin{theorem}[Accelerated hallucination inevitability]
\label{thm:accelerated-inevitability}
Assume:
\begin{enumerate}
  \item baseline context degradation \(\delta\) is monotone and
        non-invertible and satisfies
        Assumption~\ref{assump:baseline-degradation};
  \item confidence inflation \(\Phi\) with
        \(\varphi(\bar{\pi}) > 0\) for some \(\bar{\pi}\in(0,1]\);
  \item requirement escalation: \(\Req\) is monotone in credence;
  \item feedback-coupled degradation \(\delta^\star\) is credence-monotone
        and reduces to \(\delta\) at zero credence;
  \item the run contains infinitely many \emph{surface-plausible}
        explicit artifacts with \(\pi(e_n) \ge \bar{\pi}\) and recurring
        nontrivial demands relative to some floor \(r\).
\end{enumerate}
Then:
\begin{enumerate}
  \item hallucination occurs in finite time; i.e., there exists
        \(n < \infty\) such that \(t_n \prec \Req(e_n,c_n)\);
  \item the minimum such \(n\) is strictly smaller than in the baseline
        degradation model without feedback (for any fixed initial
        state and pipeline sequence), in the sense that feedback cannot
        increase time-to-hallucination.
\end{enumerate}
\end{theorem}

\begin{proof}[Proof sketch]
By unbounded credence under plausibility and the presence of infinitely
many plausible artifacts, credence \(c_n\) grows without bound along
some subsequence. By credence-monotonicity of \(\delta^\star\) and
requirement escalation, higher credence simultaneously:
\begin{itemize}
  \item increases required grounding \(\Req(e_n,c_n)\), and
  \item accelerates degradation of available grounding via
        \(\delta^\star\).
\end{itemize}
Thus the gap
\[
t_n \prec \Req(e_n,c_n)
\]
is reached in finite time, and strictly earlier than under the baseline
degradation \(\delta\) alone, where the same sequence of artifacts
would produce slower loss of grounding and lower requirements.
A fully formal proof can be obtained by coupling two runs (with and
without feedback) and using the monotone measure \(\mu\) to compare
times at which \(\mu(t_n)\) falls below \(\mu(\Req(e_n,c_n))\).
\end{proof}


\subsubsection{Degradation--repair dynamics and entropy drift}
\label{app:delta-rho}
We extend the Epistemic Dynamics to include abstract repair. The aim is to state conditions under
which repair delays or reshapes, but does not remove, long-run degradation.

\begin{definition}[Entropy functional on $T$]
An \emph{entropy functional} on the tacit category $T$ is a map $S : T \to \mathbb{R}$ such that
\begin{enumerate}
  \item $t \preceq t' \Rightarrow S(t) \le S(t')$ for all $t,t' \in T$; and
  \item there exists $S_{\mathrm{req}} \in \mathbb{R}$ such that $S(t) \le S_{\mathrm{req}}$ whenever $t$
  meets the relevant requirement floor(s), and $S(t) > S_{\mathrm{req}}$ implies that accepting a
  nontrivial demand at $t$ yields an FRFP-hallucination in the sense $t_n \prec \mathrm{Req}(e_n,c_n)$.
\end{enumerate}
Thus $S(t)$ measures a grounding deficit: larger values correspond to more degraded tacit states.
\end{definition}

We consider runs $(t_n)_{n \ge 0}$ in $T$ evolving under a mixture of degradation and repair, and
write $S_n := S(t_n)$. Let $D_n,R_n \in \{0,1\}$ indicate, respectively, whether step $n$ processes a
nontrivial demand and whether a repair operation is applied.

\begin{assumption}[Task-induced degradation]\label{ass:delta-entropy}
There exists $\varepsilon > 0$ such that for all $t \in T$ and all nontrivial demand steps to which the
Epistemic Dynamics apply,
\[
  S(\delta(t)) \;\ge\; S(t) + \varepsilon.
\]
\end{assumption}

\begin{assumption}[Bounded repair]\label{ass:rho-entropy}
There exists $B < \infty$ such that for all admissible repair operators $\rho : T \to T$ and all
$t \in T$,
\[
  S(\rho(t)) \;\ge\; S(t) - B.
\]
\end{assumption}

Combining these, for each $N$ we have
\begin{equation}\label{eq:Sn-balance}
  S_N \;\ge\; S_0 + \varepsilon \sum_{n=1}^N D_n \;-\; B \sum_{n=1}^N R_n,
\end{equation}
where the sums record entropy produced by nontrivial tasks and removed by repair.
\begin{theorem}[Deterministic entropy drift and inevitability]\label{thm:delta-rho-deterministic}
Suppose Assumptions~\ref{ass:delta-entropy} and~\ref{ass:rho-entropy} hold and that
$(D_n,R_n)_{n \ge 1}$ satisfies
\[
  \liminf_{N \to \infty} \frac{1}{N}\sum_{n=1}^N D_n \;\ge\; \alpha > 0,
  \qquad
  \limsup_{N \to \infty} \frac{1}{N}\sum_{n=1}^N R_n \;\le\; \beta < \infty.
\]
Assume further that there exists $\gamma > 0$ with
\[
  \varepsilon \alpha - B \beta \;\ge\; \gamma > 0.
\]
Then there exists $N_0$ such that $S_N \ge S_0 + \gamma N$ for all $N \ge N_0$. In particular, for
any finite $S_{\mathrm{req}}$ the hitting time
\[
  \sigma_{\mathrm{req}} := \inf\{n \ge 0 : S_n > S_{\mathrm{req}}\}
\]
is finite. If nontrivial demands recur after $\sigma_{\mathrm{req}}$ and the protocol sometimes accepts
artifacts at those steps with credence corresponding to $S_{\mathrm{req}}$, then along some admissible
run an FRFP-hallucination event $t_n \prec \mathrm{Req}(e_n,c_n)$ occurs.
\end{theorem}

\begin{proof}[Proof sketch]
Divide~\eqref{eq:Sn-balance} by $N$ and apply the liminf/limsup bounds on the empirical
frequencies of $D_n$ and $R_n$ to obtain
\[
  \liminf_{N \to \infty} \frac{S_N - S_0}{N}
  \;\ge\; \varepsilon \alpha - B \beta \;\ge\; \gamma > 0.
\]
Hence $S_N \ge S_0 + \gamma N$ for all large $N$, so $S_N$ eventually exceeds any finite
$S_{\mathrm{req}}$ and $\sigma_{\mathrm{req}} < \infty$. The hallucination clause follows from the
definition of $S_{\mathrm{req}}$ and the criterion $t_n \prec \mathrm{Req}(e_n,c_n)$ at an accepted
nontrivial demand.
\end{proof}
\begin{theorem}[Stochastic entropy drift and almost-sure inevitability]\label{thm:delta-rho-stochastic}
Suppose Assumptions~\ref{ass:delta-entropy} and~\ref{ass:rho-entropy} hold and that
$(D_n,R_n)_{n \ge 1}$ is stationary and ergodic with
\[
  \mathbb{E}[D_1] = \bar{d} > 0,
  \qquad
  \mathbb{E}[R_1] = \bar{r} < \infty.
\]
Assume that
\[
  \varepsilon \bar{d} - B \bar{r} \;=\; \gamma > 0.
\]
Then almost surely
\[
  \liminf_{N \to \infty} \frac{S_N - S_0}{N} \;\ge\; \gamma > 0,
\]
so $S_N \to +\infty$ almost surely. In particular, the hitting time
$\sigma_{\mathrm{req}} := \inf\{n \ge 0 : S_n > S_{\mathrm{req}}\}$ is almost surely finite. Under the
same acceptance and recurrence conditions as in Theorem~\ref{thm:delta-rho-deterministic}, the
hallucination time $\tau(\omega)$ in the Epistemic Dynamics stochastic semantics satisfies $P(\tau < \infty) = 1$.
\end{theorem}

\begin{proof}[Proof sketch]
By the ergodic theorem, almost surely
\[
  \lim_{N \to \infty} \frac{1}{N}\sum_{n=1}^N D_n = \bar{d},
  \qquad
  \lim_{N \to \infty} \frac{1}{N}\sum_{n=1}^N R_n = \bar{r}.
\]
Together with~\eqref{eq:Sn-balance} this yields
\[
  \liminf_{N \to \infty} \frac{S_N - S_0}{N} \;\ge\; \varepsilon \bar{d} - B \bar{r} = \gamma > 0
  \quad\text{a.s.}
\]
Thus $S_N \to +\infty$ almost surely and $\sigma_{\mathrm{req}} < \infty$ almost surely; whenever a
nontrivial demand is accepted at or after $\sigma_{\mathrm{req}}$, the hallucination criterion is met at
some finite step, so $P(\tau < \infty) = 1$.
\end{proof}
\begin{remark}[Counter-scenarios and safe horizons]\label{rem:delta-rho-counter}
The drift condition in Theorems~\ref{thm:delta-rho-deterministic} and~\ref{thm:delta-rho-stochastic}
formalizes a high-load, bounded-repair regime: nontrivial demands occur with positive asymptotic
frequency, repair has bounded one-step effect, and net entropy production $\varepsilon \bar{d}$
exceeds maximal entropy removal $B \bar{r}$. In this regime, repair changes the timing and shape
of hallucination risk but does not eliminate it.

Several counter-scenarios lie at or beyond this domain. In \emph{balanced or negative drift}
regimes, repair investment is large enough that $B \bar{r} \ge \varepsilon \bar{d}$ and $S_N$ can remain
tight or decrease. In \emph{unbounded repair} regimes, one allows repair moves with no global
bound $B$ (e.g.\ costless, arbitrarily frequent ``oracle refresh'' operations or always-perfect staff
replacement), violating Assumption~\ref{ass:rho-entropy}. In \emph{low-demand} regimes, such as
rare emergency committees, the effective demand frequency $\bar{d}$ may be so small that modest
repair keeps $S$ below $S_{\mathrm{req}}$ over relevant horizons.

Within the high-load, bounded-repair regime, improving repair capacity modifies the hazard profile
$(h_n)_{n \ge 0}$ and increases the safe-horizon functional $H_P(\varepsilon)$, but for any fixed
stochastic FRFP pipeline $P$ satisfying the drift condition, $H_P(\varepsilon)$ remains finite for every
$\varepsilon \in (0,1)$: one cannot obtain a global guarantee of ``no hallucinations ever'' over
unbounded horizons by bounded repair alone.
\end{remark}

%======================================================================
\subsubsection{Stochastic Pipelines and Their Impact}
\label{sec:B15-stochastic}
%======================================================================

We extend the preceding results to stochastic FRFP pipelines.

\subsubsubsection*{Trajectories and probability measures}

\begin{definition}[Trajectory]
Fix an initial extended state \((t_0,c_0,e_0)\in T\times
\mathbb{R}_{\ge 0}\times E\). A \emph{trajectory} (or \emph{run}) is a
sequence
\[
\omega
 =
\bigl( (t_0,c_0,e_0), (t_1,c_1,e_1), (t_2,c_2,e_2), \dots \bigr)
\]
such that each \((t_{n+1},c_{n+1},e_{n+1})\) is obtained from
\((t_n,c_n,e_n)\) by one application of the FRFP pipeline, tacit
evaluation, confidence update, and degradation (possibly
feedback-coupled).
\end{definition}

Let \(\Omega\) denote the set of all such trajectories. Let
\(\Omega_{\mathrm{allowed}}\subseteq \Omega\) be the subset of
trajectories that satisfy the structural assumptions of this appendix
(degradation, recurring requirements, etc.).

We consider two equivalent presentations of stochasticity:
\begin{itemize}
  \item \emph{Random-seed semantics:} extend the state space with a
        random seed \(r\), and suppose there is a deterministic update
        map
        \[
        F: T^e\times E \times R \to T^e \times E \times R
        \]
        such that each \(r\in R\) induces a unique trajectory
        \(\omega(r)\in\Omega\).
  \item \emph{Markov-kernel semantics:} one-step evolution is described
        by a Markov kernel
        \[
        K : T^e\times E \to \mathcal{D}(T^e\times E),
        \]
        where \(\mathcal{D}(\cdot)\) indicates probability
        distributions; a trajectory \(\omega\) is then a sample path of
        the stochastic process generated by \(K\).
\end{itemize}

In either case, a \emph{stochastic FRFP pipeline} is identified with a
probability measure \(\mathbb{P}\) on \(\Omega\) induced by the
underlying randomness.

\subsubsubsection*{Hallucination events}

Recall the hallucination criterion: hallucination at step \(n\) occurs
when
\[
t_n \prec \Req(e_n,c_n).
\]

\begin{definition}[Hallucination events]
For a trajectory \(\omega\in\Omega\) and \(n\ge 0\), define the
indicator
\[
H_n(\omega)
 :=
 \mathbf{1}\bigl[t_n \prec \Req(e_n,c_n)\bigr] \in \{0,1\}.
\]
Define the \emph{eventual hallucination indicator}
\[
H(\omega)
 :=
 \mathbf{1}\bigl[ \exists n\ge 0 : H_n(\omega)=1 \bigr].
\]
Thus \(H(\omega)=1\) iff the trajectory \(\omega\) contains at least one
hallucination.
\end{definition}

\begin{assumption}[Pathwise inevitability]
\label{assump:pathwise-inevitability}
Suppose the structural assumptions of Sections~B.1--B.12 hold (on
degradation, requirements, and feedback), and that the baseline or
accelerated hallucination inevitability theorem has been established in
the deterministic setting. We assume:
\[
\forall \omega \in \Omega_{\mathrm{allowed}}:\quad H(\omega)=1.
\]
That is, every allowed infinite trajectory eventually produces a
hallucination.
\end{assumption}

\begin{theorem}[Almost-sure inevitability under stochastic pipelines]
\label{thm:stochastic-inevitability}
Let \(\mathbb{P}\) be any probability measure on \(\Omega\) whose
support is contained in \(\Omega_{\mathrm{allowed}}\). Then
\[
\mathbb{P}\bigl( \text{eventual hallucination} \bigr)
 =
\mathbb{P}\bigl( \{\omega : H(\omega)=1\} \bigr)
 =
1.
\]
\end{theorem}

\begin{proof}
By Assumption~\ref{assump:pathwise-inevitability}, \(H(\omega)=1\) for
all \(\omega\in\Omega_{\mathrm{allowed}}\). Since \(\mathbb{P}\) is
supported on \(\Omega_{\mathrm{allowed}}\), we have \(H(\omega)=1\) for
\(\mathbb{P}\)-almost every \(\omega\). Hence
\[
\mathbb{P}(\text{eventual hallucination})
 =
\int_{\Omega} H(\omega)\, d\mathbb{P}(\omega)
 =
\int_{\Omega} 1\, d\mathbb{P}(\omega)
 =
1.
\]
\end{proof}

Thus passing from deterministic to stochastic pipelines does not weaken
the inevitability result: randomness merely selects \emph{which}
hallucinating trajectory is realized, not whether hallucination occurs.

\subsubsubsection*{Probabilistic confluence}

Stochasticity breaks single-outcome confluence but yields a natural
probabilistic analogue.

\begin{definition}[Output distribution and probabilistic confluence]
Fix a time \(n\ge 0\) and initial state \((t_0,c_0,e_0)\). Let
\(\mathbb{P}\) be a stochastic FRFP pipeline with support in
\(\Omega_{\mathrm{allowed}}\). The \emph{time-\(n\) output
distribution} is the pushforward measure
\[
\mu_n^{\mathbb{P}} := (e_n)_\# \mathbb{P},
\]
where \(e_n:\Omega\to E\) maps a trajectory to its explicit artifact at
step \(n\).

Two stochastic implementations \(\mathbb{P}_1,\mathbb{P}_2\) are
\emph{probabilistically confluent at time \(n\)} if they induce the same
output distribution \(\mu_n^{\mathbb{P}_1}=\mu_n^{\mathbb{P}_2}\) on
\(E\).
\end{definition}

If both implementations correspond to the same Markov kernel on
\(T^e\times E\), they are probabilistically confluent at every time
\(n\). This notion provides a distributional counterpart to the
canonical-form and confluence properties discussed earlier.

%======================================================================
\subsubsection{Finite-Horizon Pipelines and Their Impact}
%======================================================================

We now study the effect of restricting attention to finite prefixes of
trajectories, both in deterministic and stochastic settings.

\subsubsubsection*{Hallucination time}

\begin{definition}[Hallucination time]
For a trajectory \(\omega\in\Omega\), define the \emph{hallucination
time}
\[
\tau(\omega)
 :=
 \inf\{ n \ge 0 : t_n \prec \Req(e_n,c_n) \},
\]
with the convention \(\tau(\omega)=+\infty\) if the set is empty (i.e.,
if no hallucination ever occurs).
\end{definition}

Thus \(\tau(\omega) < \infty\) iff \(\omega\) contains at least one
hallucination. In terms of Definition~B.13.2, \(H(\omega) =
\mathbf{1}[\tau(\omega) < \infty]\).

\begin{lemma}[Pathwise inevitability and finite hallucination time]
Under Assumption~\ref{assump:pathwise-inevitability} (pathwise
inevitability), we have
\[
\forall \omega \in \Omega_{\mathrm{allowed}}:\quad \tau(\omega) < \infty.
\]
\end{lemma}

\begin{proof}
By Assumption~\ref{assump:pathwise-inevitability}, for each
\(\omega\in\Omega_{\mathrm{allowed}}\) there exists some \(n\) with
\(H_n(\omega)=1\). By definition of \(\tau\), this implies that
\(\tau(\omega)\) is the infimum of a nonempty subset of \(\mathbb{N}\),
hence finite.
\end{proof}

A \emph{finite horizon} of length \(N\) corresponds to observing only
the prefix
\[
\omega_{\le N}
 :=
((t_0,c_0,e_0),\dots,(t_N,c_N,e_N)).
\]
We say that a trajectory \emph{hallucinates within horizon \(N\)} if
\(\tau(\omega)\le N\).

\begin{remark}[No universal finite bound in general]
Pathwise inevitability (\(\tau(\omega) < \infty\) for all allowed
\(\omega\)) does not imply the existence of a universal finite upper
bound \(N^\ast\) such that \(\tau(\omega)\le N^\ast\) for all
\(\omega\in\Omega_{\mathrm{allowed}}\). In general, for any fixed
\(N\) there may exist trajectories with \(\tau(\omega) > N\).
\end{remark}

\subsubsubsection*{A finite-time bound under strengthened assumptions}

If desired, one can obtain a finite-time upper bound under stronger
assumptions.

\begin{assumption}[Uniform loss and requirement floor]
\label{assump:uniform-loss}
Suppose there exist:
\begin{itemize}
  \item a grounding measure \(\mu:T\to\mathbb{R}_{\ge 0}\);
  \item a requirement level \(r\in T\) and constant \(\alpha>0\) with
        \(\mu(r)=\alpha\);
  \item a constant \(\varepsilon>0\);
\end{itemize}
such that, for every allowed trajectory
\(\omega=((t_n,c_n,e_n))_{n\ge 0}\):
\begin{enumerate}
  \item \emph{Monotone loss:} \(\mu(t_{n+1})\le\mu(t_n)\) for all
        \(n\).
  \item \emph{Uniform loss above \(r\):} whenever \(\mu(t_n)>\alpha\),
        we have
        \[
        \mu(t_n) - \mu(t_{n+1}) \ge \varepsilon.
        \]
  \item \emph{Recurrent requirement:} there exists an infinite sequence
        \((n_k)\) with \(\Req(e_{n_k},c_{n_k})\succeq r\) for all
        \(k\).
\end{enumerate}
\end{assumption}

\begin{lemma}[Finite-time bound under uniform loss]
\label{lem:finite-time-bound}
Under Assumption~\ref{assump:uniform-loss}, every allowed trajectory
\(\omega\) satisfies
\[
\tau(\omega)
 \le
N^\ast
 :=
\left\lceil \frac{\mu(t_0) - \alpha}{\varepsilon} \right\rceil.
\]
That is, a hallucination must occur by time \(N^\ast\).
\end{lemma}

\begin{proof}
By monotone loss and the uniform-loss condition, as long as
\(\mu(t_n)>\alpha\),
\[
\mu(t_{n+1}) \le \mu(t_n) - \varepsilon.
\]
By induction, for any \(n\) with this property,
\[
\mu(t_n) \le \mu(t_0) - n\varepsilon.
\]
For \(n = N^\ast\) we obtain \(\mu(t_{N^\ast}) \le \alpha =
\mu(r)\). By recurrence of requirements, there exists some \(k\) with
\(n_k \le N^\ast\) and \(\Req(e_{n_k},c_{n_k}) \succeq r\). Using
monotonicity of \(\mu\) with respect to \(\preceq\), we conclude that
\[
t_{n_k} \preceq r \preceq \Req(e_{n_k},c_{n_k}),
\]
and in fact \(t_{n_k} \prec \Req(e_{n_k},c_{n_k})\) for nontrivial
requirements, so \(\tau(\omega)\le n_k \le N^\ast\).
\end{proof}

These stronger conditions are not needed for the main inevitability
results, but they demonstrate that finite-time bounds are available in
regimes with sufficiently uniform context loss and recurrent demands.

\subsubsubsection*{Stochastic finite horizons}

Returning to the stochastic setting, we can describe how
finite-horizon hallucination probabilities behave.

\begin{definition}[Finite-horizon hallucination probability]
Fix a horizon \(N\in\mathbb{N}\). Define the event
\[
A_N := \{\omega\in\Omega : \tau(\omega) \le N\},
\]
and the corresponding probability
\[
p_N := \mathbb{P}(A_N)
     = \mathbb{P}(\tau \le N).
\]
Thus \(p_N\) is the probability that a random run hallucinates within
the first \(N\) steps.
\end{definition}

\begin{proposition}[Monotone convergence of finite-horizon probabilities]
Assume almost-sure inevitability
(Theorem~\ref{thm:stochastic-inevitability}), i.e.,
\(\mathbb{P}(\tau<\infty)=1\). Then:
\begin{enumerate}
  \item The sequence \((p_N)_{N\ge 0}\) is non-decreasing:
        \(p_N \le p_{N+1}\) for all \(N\).
  \item \(\displaystyle \lim_{N\to\infty} p_N = 1\).
\end{enumerate}
\end{proposition}

\begin{proof}
For each \(\omega\), the events \(\{\tau(\omega)\le N\}\) are nested:
\(\{\tau\le N\}\subseteq\{\tau\le N+1\}\). Hence
\(A_N\subseteq A_{N+1}\) and \(p_N\le p_{N+1}\). Moreover,
\(\{\tau<\infty\}=\bigcup_{N} A_N\), so by continuity of probability
from below,
\[
\mathbb{P}(\tau<\infty)
 =
\lim_{N\to\infty} \mathbb{P}(A_N)
 =
\lim_{N\to\infty} p_N.
\]
By assumption, \(\mathbb{P}(\tau<\infty)=1\), so
\(\lim_{N\to\infty} p_N=1\).
\end{proof}

\begin{proposition}[Multiple runs over a fixed horizon]
Fix a horizon \(N\) and consider \(K\in\mathbb{N}\) independent
trajectories \(\omega^{(1)},\dots,\omega^{(K)}\), with corresponding
hallucination times \(\tau^{(i)}\). Then:
\begin{enumerate}
  \item The probability that no trajectory hallucinates within horizon
        \(N\) is
        \[
        \mathbb{P}\bigl(\forall i,\ \tau^{(i)} > N\bigr)
         = (1-p_N)^K.
        \]
  \item The probability that at least one trajectory hallucinates
        within horizon \(N\) is
        \[
        1 - (1-p_N)^K.
        \]
\end{enumerate}
\end{proposition}

\begin{proof}
By independence,
\[
\mathbb{P}(\tau^{(i)} > N) = 1 - p_N
\]
for each \(i\), so
\[
\mathbb{P}\bigl(\forall i,\ \tau^{(i)} > N\bigr)
 =
\prod_{i=1}^K (1-p_N)
 =
(1-p_N)^K.
\]
The complementary event has probability \(1-(1-p_N)^K\).
\end{proof}

These results show how inevitability manifests at finite horizons: for
any fixed horizon, hallucination remains probabilistic (\(p_N<1\) in
general), but as the horizon length or number of runs increases, the
probability that hallucination has occurred approaches~1.

\medskip

Future work lifts these individual dynamics to populations of agents and
explicit consensus mechanisms.
% ------------------------------------------------------------------
% New material after Proposition B.38 in Appendix B
% ------------------------------------------------------------------

\subsubsection*{B.16 Hazard Functions and Safe Horizons}

We now refine finite-horizon analysis by introducing per-step hazard and a derived notion of
``safe horizon'' for a given stochastic FRFP pipeline.

\begin{definition}[Per-step hazard]\label{def:hazard}
Let $\tau : \Omega \to \mathbb{N} \cup \{\infty\}$ be the hallucination time of Definition~B.31,
and let $P$ be a stochastic FRFP pipeline with associated distribution of $\tau$.
For each $n \in \mathbb{N}$ with $P(\tau \ge n) > 0$ we define the (discrete) hazard at step $n$ by
\[
  h_n := P(\tau = n \mid \tau \ge n)
       = \frac{P(\tau = n)}{P(\tau \ge n)}.
\]
If $P(\tau \ge n) = 0$ we set $h_n := 0$.
\end{definition}

Intuitively, $h_n$ is the conditional probability that a run first hallucinates at step $n$,
given that it has not hallucinated earlier. The sequence $(h_n)_{n \ge 0}$ is a pipeline-level
summary of how risk is distributed along the trajectory.

\begin{definition}[Finite-horizon survival and hazard]\label{def:survival}
For $N \in \mathbb{N}$ let
\[
  S_N := P(\tau > N) = 1 - p_N
\]
be the survival probability at horizon $N$, where $p_N = P(\tau \le N)$ is the finite-horizon
hallucination probability from Definition~B.36. Whenever $P(\tau \ge n) > 0$ for all $n$,
the survival probabilities admit the product representation
\[
  S_N = \prod_{n=0}^{N} (1 - h_n),
\]
so that the finite-horizon risk $p_N = 1 - S_N$ can be understood as the cumulative effect
of per-step hazards.
\end{definition}

This representation makes explicit the trade-off between (i) the number of steps in a pipeline
and (ii) the hazard profile $(h_n)$ induced by its design (e.g. decomposition, tool calls,
or repeated self-checks).

\begin{definition}[Safe horizon at tolerance $\varepsilon$]\label{def:safe-horizon}
Fix a tolerance level $\varepsilon \in (0,1)$. The $\varepsilon$-safe horizon of a stochastic
FRFP pipeline $P$ is defined as
\[
  H_P(\varepsilon)
  := \sup\{\,N \in \mathbb{N} : p_N \le \varepsilon\,\}
  \in \mathbb{N} \cup \{\infty\}.
\]
When $H_P(\varepsilon)$ is finite, any horizon $N \le H_P(\varepsilon)$ satisfies
$P(\tau \le N) \le \varepsilon$, i.e. hallucination within that horizon is at most
$\varepsilon$-likely.
\end{definition}

\begin{remark}[Design interpretation]
The safe-horizon functional $H_P(\varepsilon)$ isolates a pipeline-level design parameter:
for a given stochastic implementation $P$ and admissible risk budget $\varepsilon$, it
characterizes how long a single run can safely be extended before hallucination becomes
probable. Changes to decomposition, step-ordering, or intermediate checks can therefore
be compared in terms of their effect on $(h_n)$ and $H_P(\varepsilon)$, even when the
underlying explicit pipeline is statically equivalent in $N \ltimes E$.
\end{remark}

Combining Definition~\ref{def:safe-horizon} with Proposition~B.38, we can also analyze
restart and parallelization policies: for example, running $K$ independent instances of a
pipeline up to horizon $N \le H_P(\varepsilon)$ yields hallucination probability
$1 - (1 - p_N)^K$, which can be bounded using the same tolerance parameter.

\subsubsection*{B.17 Dynamic Refinement of Pipelines}

We next introduce a notion of dynamic refinement between stochastic FRFP pipelines that
compares their hallucination-time distributions.

\begin{definition}[Dynamic refinement]\label{def:dynamic-refinement}
Let $P_1, P_2$ be stochastic FRFP pipelines on the same state space and admissible trajectory
set $\Omega_{\mathrm{allowed}}$, with induced hallucination times $\tau_1, \tau_2$ and
finite-horizon probabilities
\[
  p_N^{(i)} := P_i(\tau_i \le N), \quad N \in \mathbb{N}, \; i \in \{1,2\}.
\]
We say that $P_2$ is a \emph{dynamic refinement} (or \emph{dynamically safer variant}) of $P_1$
if for all $N \in \mathbb{N}$ we have
\[
  p_N^{(2)} \;\le\; p_N^{(1)}.
\]
Equivalently, $\tau_2$ stochastically dominates $\tau_1$ in the usual sense.
\end{definition}

\begin{remark}[Static vs dynamic equivalence]
Two pipelines may be statically equivalent in the sense of Appendix~A (same normal form in
$N \ltimes E$ and same tacit semantics $J{-}K$), yet dynamically inequivalent: they can induce
different hazard profiles and hence different hallucination-time distributions.
Dynamic refinement isolates this difference at the level of trajectories, not static syntax:
$P_2$ improves on $P_1$ precisely when it never increases finite-horizon hallucination risk.
\end{remark}

\begin{remark}[Limitations]
Dynamic refinement is a partial order on stochastic implementations, not an absolute
guarantee of safety: if the baseline pipeline $P_1$ satisfies the inevitability conditions
of Theorems~B.11, B.25, and B.29, then so does any refinement $P_2$. Refinements can
delay hallucination and reduce risk over practically relevant horizons, but they cannot
eliminate hallucination altogether without violating the structural assumptions of this
appendix.
\end{remark}
% ------------------------------------------------------------------
% New subsection at the end of Appendix B
% ------------------------------------------------------------------

\subsubsection*{B.18 Structural implications of pipeline dynamics}

The move from static FRFP pipelines to the dynamic setting of this appendix
turns time and iteration into first-class objects. Beyond the inevitability
results of Theorems~B.11, B.25, and B.29, this buys several structural
notions that will be used in later appendices and case studies.

\paragraph{(1) Hazard profiles and safe horizons.}
Definition~B.31 endows each stochastic FRFP implementation with a
hallucination time $\tau$, and Definitions~B.36 and~\ref{def:hazard} refine
this to a finite-horizon risk curve $(p_N)_{N \ge 0}$ and a discrete hazard
sequence $(h_n)_{n \ge 0}$. The safe-horizon functional
$H_P(\varepsilon)$ in Definition~\ref{def:safe-horizon} abstracts how long a
given pipeline can be run, under a fixed stochastic implementation $P$ and
risk budget $\varepsilon$, before hallucination becomes likely. Thus even
when static semantics and correctness are fixed, pipeline dynamics exposes a
new design degree of freedom: the shape of the hazard profile along a run.

\paragraph{(2) Time--decomposition trade-offs.}
In the static theory of Appendices~A and~B.1--B.3, different decompositions
of a task into explicit steps are equivalent whenever they normalize to the
same morphism in $N \ltimes E$ and share the same tacit semantics $J{-}K$.
Once trajectories are considered, this equivalence is refined: more
fine-grained decompositions typically reduce the requirement demanded at
individual steps, but also increase the length of the trajectory and the
opportunity for degradation to act. Pipeline dynamics therefore makes precise
a basic design tension: reshaping a task into ``more steps'' can both reduce
per-step hazard and increase cumulative risk, and only the dynamic
quantities $(h_n)$ and $(p_N)$ decide which effect dominates in a concrete
setting.

\paragraph{(3) Dynamic refinement of implementations.}
Definition~\ref{def:dynamic-refinement} introduces a partial order on
stochastic implementations of the \emph{same} static pipeline: $P_2$
dynamically refines $P_1$ when it never increases finite-horizon risk.
This separates:
\begin{itemize}
  \item static equivalence in $N \ltimes E$ (same normal form, same tacit
        semantics), from
  \item dynamic safety properties (how quickly hallucination emerges under a
        particular scheduling, tool use pattern, or feedback coupling).
\end{itemize}
Subsequent appendices will use this distinction to talk about pipeline-level
improvements that do not change what a system computes in principle, but do
change how its risk unfolds in time.

\paragraph{(4) Session limits, restart policies, and composition.}
Because hallucination is represented as a stopping time and finite-horizon
probabilities are explicit, Appendix~B supports principled discussion of
session limits and restarts. Proposition~B.38 already shows how repeated
runs over a fixed horizon aggregate risk via $1 - (1 - p_N)^K$. In more
complex settings, this allows single-pipeline guarantees (bounds on $p_N$
or $H_P(\varepsilon)$) to be composed into system-level guarantees for
restart, parallelization, or human-in-the-loop oversight policies, much as
Appendix~A.10 composes agent-level inevitability into collective inevitability.

\paragraph{(5) Localizing and targeting failure modes.}
Finally, working with trajectories rather than only normal forms means that
hallucination can be localized to segments of a run: particular regions of
the state space, update rules, or phases of the pipeline where hazard is
concentrated. This enables targeted interventions (for example, inserting
human checks or additional tools only before entering high-hazard regions)
that would be invisible in a purely static semantics.

Taken together, these structures show that pipeline dynamics is not only a
way to restate hallucination inevitability in stochastic language. It also
provides a vocabulary---hazard, safe horizons, dynamic refinement, and
time--decomposition trade-offs---for comparing and improving FRFP pipelines
\emph{without} altering their underlying explicit semantics in $N \ltimes E$.
